% GENERATED FILE. Edit canonical lessons, then run npm run generate:books.
% canonical-source-hash: fnv1a64:81d218a4a49ccc50
% canonical-lessons: UR-C230-poti, UR-C230-sasur, UR-C230-sas, UR-C230-dulha, UR-C230-dulhan

\chapter{Granddaughter, Father-in-law, Mother-in-law, Bridegroom, Bride}
\label{ch:ur-granddaughter-father-in-law-mother-in-law-bridegroom-bride}
\hlchaptermodality{230}

\begin{chapteropening}
\textbf{By the end of this chapter:} \emph{I can say a granddaughter, a father-in-law, a mother-in-law, a bridegroom and a bride.}

Complete the last lesson of chapter 230: I can say a granddaughter, a father-in-law, a mother-in-law, a bridegroom and a bride.
\end{chapteropening}

\section[\texorpdfstring{\ur{پوتی} (potī)}{پوتی (potī)}]{\ur{پوتی} (potī) — a granddaughter (son's daughter)}

\label{lesson:UR-C230-poti}



\begin{quote}
Before the new one: say the Urdu for a petrol station, then the Urdu for a seat.
\end{quote}

\subsection*{You'll want to know: \ur{پوتی}}
\textbf{\ur{پوتی}} — \emph{potī} — \textquotedblleft{}a granddaughter (son's daughter)\textquotedblright{}.

\begin{cousinweb}
\textbf{\ur{پوتی}} (\emph{potī}) is a tadbhava, a word inherited from Sanskrit and reshaped by sound change. It comes from \emph{pautrī}, \textquotedblleft{}son's daughter\textquotedblright{} (CDIAL 8417). It is the feminine of \emph{pautra}, \textquotedblleft{}son's son\textquotedblright{}.
\end{cousinweb}

\begin{grammarlens}[title={What you've built}]
One more word for this chapter.
\end{grammarlens}

\subsection*{Your turn}
\begin{itemize}
\raggedright
  \item \emph{Say it:} \emph{potī}
  \item \emph{Say it:} \emph{potī}, once more
  \item \emph{Say it:} say \emph{potī}
  \item \emph{Recall:} say the Urdu for a petrol station, then the Urdu for a seat, then say \emph{potī} again
\end{itemize}

\begin{tcolorbox}[breakable,colback=teal!4,colframe=teal!35!black,title={Before you move on}]
What does \ur{پوتی} mean? (\textquotedblleft{}A granddaughter (son's daughter)\textquotedblright{}.) Say it once more.
\end{tcolorbox}
\section[\texorpdfstring{\ur{سسر} (sasur)}{سسر (sasur)}]{\ur{سسر} (sasur) — a father-in-law}

\label{lesson:UR-C230-sasur}



\begin{quote}
Before the new one: say the Urdu for a seat, then the Urdu for a granddaughter.
\end{quote}

\subsection*{You'll want to know: \ur{سسر}}
\textbf{\ur{سسر}} — \emph{sasur} — \textquotedblleft{}a father-in-law\textquotedblright{}.

\begin{cousinweb}
\textbf{\ur{سسر}} (\emph{sasur}) is a tadbhava, a word inherited from Sanskrit and reshaped by sound change. It comes from \emph{śvaśura}, \textquotedblleft{}father-in-law\textquotedblright{} (CDIAL 12753).
\end{cousinweb}

\begin{grammarlens}[title={What you've built}]
One more word for this chapter.
\end{grammarlens}

\subsection*{Your turn}
\begin{itemize}
\raggedright
  \item \emph{Say it:} \emph{sasur}
  \item \emph{Say it:} \emph{sasur}, once more
  \item \emph{Say it:} say \emph{sasur}
  \item \emph{Recall:} say the Urdu for a seat, then the Urdu for a granddaughter, then say \emph{sasur} again
\end{itemize}

\begin{tcolorbox}[breakable,colback=teal!4,colframe=teal!35!black,title={Before you move on}]
What does \ur{سسر} mean? (\textquotedblleft{}A father-in-law\textquotedblright{}.) Say it once more.
\end{tcolorbox}
\section[\texorpdfstring{\ur{ساس} (sās)}{ساس (sās)}]{\ur{ساس} (sās) — a mother-in-law}

\label{lesson:UR-C230-sas}



\begin{quote}
Before the new one: say the Urdu for a granddaughter, then the Urdu for a father-in-law.
\end{quote}

\subsection*{You'll want to know: \ur{ساس}}
\textbf{\ur{ساس}} — \emph{sās} — \textquotedblleft{}a mother-in-law\textquotedblright{}.

\begin{cousinweb}
\textbf{\ur{ساس}} (\emph{sās}) is a tadbhava, a word inherited from Sanskrit and reshaped by sound change. It comes from \emph{śvaśrū}, \textquotedblleft{}husband's or wife's mother\textquotedblright{} (CDIAL 12759).
\end{cousinweb}

\begin{grammarlens}[title={What you've built}]
One more word for this chapter.
\end{grammarlens}

\subsection*{Your turn}
\begin{itemize}
\raggedright
  \item \emph{Say it:} \emph{sās}
  \item \emph{Say it:} \emph{sās}, once more
  \item \emph{Say it:} say \emph{sās}
  \item \emph{Recall:} say the Urdu for a granddaughter, then the Urdu for a father-in-law, then say \emph{sās} again
\end{itemize}

\begin{tcolorbox}[breakable,colback=teal!4,colframe=teal!35!black,title={Before you move on}]
What does \ur{ساس} mean? (\textquotedblleft{}A mother-in-law\textquotedblright{}.) Say it once more.
\end{tcolorbox}
\section[\texorpdfstring{\ur{دولہا} (dūlhā)}{دولہا (dūlhā)}]{\ur{دولہا} (dūlhā) — a bridegroom}

\label{lesson:UR-C230-dulha}



\begin{quote}
Before the new one: say the Urdu for a father-in-law, then the Urdu for a mother-in-law.
\end{quote}

\subsection*{You'll want to know: \ur{دولہا}}
\textbf{\ur{دولہا}} — \emph{dūlhā} — \textquotedblleft{}a bridegroom\textquotedblright{}.

\begin{cousinweb}
\textbf{\ur{دولہا}} (\emph{dūlhā}) is a tadbhava, a word inherited from Sanskrit and reshaped by sound change. It comes from \emph{durlabha}, \textquotedblleft{}hard to obtain; beloved\textquotedblright{} (CDIAL 6446).
\end{cousinweb}

\begin{grammarlens}[title={What you've built}]
One more word for this chapter.
\end{grammarlens}

\subsection*{Your turn}
\begin{itemize}
\raggedright
  \item \emph{Say it:} \emph{dūlhā}
  \item \emph{Say it:} \emph{dūlhā}, once more
  \item \emph{Say it:} say \emph{dūlhā}
  \item \emph{Recall:} say the Urdu for a father-in-law, then the Urdu for a mother-in-law, then say \emph{dūlhā} again
\end{itemize}

\begin{tcolorbox}[breakable,colback=teal!4,colframe=teal!35!black,title={Before you move on}]
What does \ur{دولہا} mean? (\textquotedblleft{}A bridegroom\textquotedblright{}.) Say it once more.
\end{tcolorbox}
\section[\texorpdfstring{\ur{دلہن} (dulhan)}{دلہن (dulhan)}]{\ur{دلہن} (dulhan) — a bride}

\label{lesson:UR-C230-dulhan}



\begin{quote}
Before the new one: say the Urdu for a mother-in-law, then the Urdu for a bridegroom.
\end{quote}

\subsection*{You'll want to know: \ur{دلہن}}
\textbf{\ur{دلہن}} — \emph{dulhan} — \textquotedblleft{}a bride\textquotedblright{}.

\begin{cousinweb}
\textbf{\ur{دلہن}} (\emph{dulhan}) is the feminine of \emph{dūlhā}, \textquotedblleft{}bridegroom\textquotedblright{}. CDIAL lists both under Sanskrit \emph{durlabha}, \textquotedblleft{}hard to obtain; beloved\textquotedblright{} (CDIAL 6446).
\end{cousinweb}

\begin{grammarlens}[title={What you've built}]
One more word for this chapter.
\end{grammarlens}

\subsection*{Your turn}
\begin{itemize}
\raggedright
  \item \emph{Say it:} \emph{dulhan}
  \item \emph{Say it:} \emph{dulhan}, once more
  \item \emph{Say it:} say \emph{dulhan}
  \item \emph{Recall:} say the Urdu for a mother-in-law, then the Urdu for a bridegroom, then say \emph{dulhan} again
\end{itemize}

\begin{tcolorbox}[breakable,colback=teal!4,colframe=teal!35!black,title={Before you move on}]
What does \ur{دلہن} mean? (\textquotedblleft{}A bride\textquotedblright{}.) Say it once more.
\end{tcolorbox}
